\section[General spectral gap certificates]{A computable spectral gap certificate and paid upward rounding}
\label{sec:op3-spectral-vwf-certificate}

\paragraph{Scope.}
The following elementary spectral/convexity comparison gives a useful
stopping interface for the supplied numerical recursion. The comparison
idea is standard; this note supplies its exact mixed-error certificate
and charged VWF realization. These are \statusproved\ drafts awaiting
independent review. They do not provide the preconditioner or its oracle.
In particular, they do not imply a local OP3 work bound.

\begin{theorem}[A nongrounded spectral residual certificate; \statusproved]
\label{thm:op3-spectral-vwf-gap-certificate}
Let $G\preceq H\preceq\kappa G$, $\kappa\geq1$, and let $a$ be
any feasible point of $\Phi(x)=x^TGx/2+h(x)$ with finite infimum
$\Phi^*$. The proper convex function $h$ includes the domain indicator.
Define
\[
 R_G(z)=\Phi(a+z)-\Phi(a)=\tfrac12z^TGz+h_a(z),\qquad
 R_H(z)=\tfrac12z^THz+h_a(z).
\]
Thus $h_a(0)=0$ and its domain is convex and contains zero.
Suppose an oracle returns feasible $q$ with
$R_H(q)\leq R_H^*/\beta+\xi$, where $\beta\geq1$ and $\xi\geq0$.
Then the computable quantity
\begin{equation}
 \mathcal C(a)=\kappa\beta\bigl(\xi-R_H(q)\bigr)
 \label{eq:op3-spectral-vwf-gap-certificate}
\end{equation}
satisfies
\[
 0\leq\Phi(a)-\Phi^*\leq\mathcal C(a)
 \leq\kappa\beta\bigl(\Phi(a)-\Phi^*+\xi\bigr).
\]
No positive grounding, objective-gap margin or coordinate box is required.
The point $a$ need not have nonpositive objective value.
\end{theorem}
\begin{proof}
The spectral order and convexity at the feasible origin give, for all
feasible $z$,
\[
 R_G(z)\leq R_H(z),\qquad
 R_H(z/\kappa)\leq R_G(z)/\kappa.
\]
The second inequality uses $H/\kappa^2\preceq G/\kappa$ and
$h_a(z/\kappa)\leq h_a(z)/\kappa$. Taking infima proves
\[
 \kappa R_H^*\leq R_G^*\leq R_H^*\leq0,
 \qquad R_G^*=-(\Phi(a)-\Phi^*).
\]
The mixed oracle contract yields
$\kappa\beta(R_H(q)-\xi)\leq\kappa R_H^*\leq R_G^*$,
which is the lower certificate inequality. Feasibility yields
$R_H(q)\geq R_H^*\geq R_G^*$, proving the upper one. The argument
uses only finite infima; attainment is unnecessary for the comparison.
\end{proof}

\begin{corollary}[Computable absolute and mixed stopping rules; \statusproved]
\label{cor:op3-spectral-vwf-stopping}
For a requested absolute gap $t>0$, accept if $\mathcal C(a)\leq t$.
If $\xi\leq t/(2\kappa\beta)$, this must accept whenever the true
gap is at most $t/(2\kappa\beta)$.
For a requested $(2,\eta)$ output, it suffices to check
\[
 \Phi(a)+\mathcal C(a)\leq2\eta.
\]
For a feasible-origin normalized objective with nonpositive infimum,
$\mathcal C(a)\leq\zeta$ also certifies its $(1+\delta,\zeta)$
contract for every $\delta>0$.
\end{corollary}
\begin{proof}
The first claim follows from both certificate inequalities.
For the second,
$2\Phi(a)\leq\Phi(a)-\mathcal C(a)+2\eta\leq\Phi^*+2\eta$.
For the third, $\Phi(a)\leq\Phi^*+\zeta
\leq\Phi^*/(1+\delta)+\zeta$.
\end{proof}

\paragraph{Charged VWF implementation.}
With original vertex functions $f_i$, the residual functions are
$f_i(a_i+z)-f_i(a_i)+(Ga)_iz$ on $[L_i-a_i,\infty)$.
They are constructed with a complete sparse product, anchor queries and
piece copies, and the oracle energy is evaluated on $H$, not on $G$.
This costs $O(m_G+m_H+S)$ work and $O(S+n)$ new words besides the
oracle's own complete costs, where $S$ is the total piece count.
Each failed certificate remains charged. When $H$ is a supplied tree,
\cref{thm:op3-persistent-vwf-forest} computes $R_H^*$ exactly using a
single retained root. Only its minimizing scalar value is needed:
no full-coordinate reconstruction is required by the certificate.
This realizes $\beta=1$, $\xi=0$, with the stated persistent-tree work
and all-allocation bound. Dense minima in the audit are validators only.

\begin{lemma}[An explicit upward rounding budget; \statusproved]
\label{lem:op3-upward-vwf-rounding}
Consider a supplied graph VWF objective with total edge weight $W$,
total largest curvature $C$, and $A=\sum_i|f_i'(0)|$. Let feasible
$x$ satisfy $\|x\|_\infty\leq R$. For a vector $d$ with
$0\leq d_i\leq s\leq1$, the point $x+d$ remains feasible and
\[
 \Phi(x+d)-\Phi(x)
 \leq s\bigl(A+(4W+C)R\bigr)+\tfrac12(W+C)s^2
 \leq s\bigl(A+(4W+C)(R+1)\bigr).
\]
Thus upward rounding to a dyadic grid with spacing at most
$\min\{1,\zeta/[2(A+(4W+C)(R+1))]\}$ raises energy by at most
$\zeta/2$. If the denominator is zero, the objective is constant and
there is no rounding error.
All range scans, grid selection, arithmetic and output are charged.
\end{lemma}
\begin{proof}
The graph gradient satisfies $\|Gx\|_1\leq4WR$, and bounded scalar
curvature gives $\sum_i|f_i'(x_i)|\leq A+CR$.
Consequently $\nabla\Phi(x)^Td\leq s(A+(4W+C)R)$.
The exact graph quadratic remainder is at most $Ws^2/2$, because
$|d_i-d_j|\leq s$. Integrating each scalar derivative gives total
vertex remainder at most $Cs^2/2$. These arguments cross all scalar
breakpoints and include lower endpoints via their right derivatives.
The displayed bound and feasibility follow. A dyadic spacing can be
found by repeated halving; this costs logarithmically many comparisons
in the displayed range and requested tolerance. Implementing integer
rounding in the rational validator is separate from a bit-complexity
claim, which is not made here.
\end{proof}

\paragraph{Relation to the nested audit.}
The fixed-length nested audit runs the mixed accelerated updates at two
levels. Each outer inner problem is eliminated once; its coarse residuals
are compressed and solved by another accelerated invocation. That
invocation uses a supplied tree preconditioner and exact persistent
single-root minima. Upward dyadic rounding deliberately makes its base
oracle inexact, and the mixed energy inequality is checked exactly.
Every refinement call, rejected candidate, normalized model, source
update, curve version and recovered coordinate is counted. Dense minima
are used only to verify gaps, never to choose a stopping iteration or
an accelerated output. A separately implemented certified variant stops
only when its paid tree certificate proves the required mixed accuracy.
For the same two physical profiles it uses 6,525 inner accelerated steps
and 12,605 additional tree minimum-value certificates, versus 389,120
inner steps for the fixed-length variant. A third signed profile with
breakpoints at $\pm2^{-100}$ exercises positive compression errors and
actual piece changes. These are exact finite integration audits, not
asymptotic or local-algorithm timing claims.
